\documentclass[11pt,a4paper]{article}
\usepackage[T1]{fontenc}
\usepackage{lmodern}
\usepackage[margin=27mm,headheight=14pt]{geometry}
\usepackage{amsmath,amssymb,amsthm,mathtools}
\usepackage{microtype,booktabs,array,longtable,enumitem}
\usepackage[dvipsnames]{xcolor}
\usepackage{hyperref}
\usepackage{fancyhdr}
\usepackage{listings}
\hypersetup{colorlinks=true,linkcolor=MidnightBlue,citecolor=MidnightBlue,urlcolor=MidnightBlue,
 pdftitle={The Lower Critical Endpoint of Exponential-Feedback Transseries},
 pdfauthor={Research manuscript prepared for Vladimir Reshetnikov}}
\pagestyle{fancy}\fancyhf{}
\fancyhead[L]{\small The lower critical endpoint}
\fancyhead[R]{\small ProveIt research continuation}
\fancyfoot[C]{\thepage}
\setlength{\emergencystretch}{2em}
\setlength{\parskip}{3pt}
\setlist[enumerate]{leftmargin=*,itemsep=4pt}
\newtheorem{theorem}{Theorem}[section]
\newtheorem{proposition}[theorem]{Proposition}
\newtheorem{lemma}[theorem]{Lemma}
\newtheorem{corollary}[theorem]{Corollary}
\theoremstyle{definition}\newtheorem{definition}[theorem]{Definition}
\theoremstyle{remark}\newtheorem{remark}[theorem]{Remark}
\newcommand{\eps}{\varepsilon}
\newcommand{\Li}{\operatorname{Li}}
\newcommand{\Prob}{\mathbb P}
\newcommand{\E}{\mathbb E}
\newcommand{\N}{\mathbb N}
\newcommand{\R}{\mathbb R}
\newcommand{\C}{\mathbb C}
\newcommand{\TV}{d_{\mathrm{TV}}}
\newcommand{\Law}{\mathcal L}
\newcommand{\Poi}{\operatorname{Pois}}
\newcommand{\1}{\mathbf 1}
\newcommand{\dd}{\,\mathrm d}
\newcommand{\coeff}[2]{[#1^{#2}]}
\newcommand{\repo}{\texttt{ProveIt}}
\newcommand{\rzero}{\rho_{\eps}}
\makeatletter
\newcommand{\inputtable}[1]{\@@input #1 }
\makeatother
\DeclareMathOperator{\supp}{supp}
\DeclareMathOperator{\ImPart}{Im}
\lstset{basicstyle=\ttfamily\small,breaklines=true,columns=fullflexible,frame=single}

\begin{document}
\begin{titlepage}
\vspace*{9mm}
\noindent\textsc{Research continuation \; / \; Analysis and probability}
\vspace{12mm}
\begin{center}
{\LARGE\bfseries The Lower Critical Endpoint of\\[5pt]
Exponential-Feedback Transseries}\par
\vspace{8mm}
{\large Uniform coefficient asymptotics, deletion of the largest action,\\[3pt]
Landau cutoff profiles, and the compound-Poisson boundary}\par
\vspace{11mm}
{Prepared for Vladimir Reshetnikov}\par
\vspace{3mm}
{29 September 2026}
\end{center}
\vspace{9mm}
\begin{abstract}
We study the lower endpoint of the critical power-law feedback equation
$U(q)=A_{\eps}(q e^{U(q)})$, where
$A_{\eps}(z)=(\Li_{2+\eps}(z)+P(z))/(\zeta(1+\eps)+P'(1))$,
$\eps>0$, and $P$ changes a fixed finite prefix while preserving
nonnegative coefficients. This is the $\alpha\downarrow1$ half of
Question~4 in the repository article \emph{Beyond Finite-Action Folds}.
Writing $\alpha=1+\eps$, $\beta=1/\alpha$,
$K_{\eps}=\Gamma(-\alpha)/(\zeta(\alpha)+P'(1))$,
$B=(nK_{\eps})^{\beta}$, and
$\rho_{\eps}=e^{-A_{\eps}(1)}$, we prove, uniformly down to arbitrarily
small positive $\eps$,
\[
 [q^n]U(q)\sim
 \frac{\rho_{\eps}^{-n}}{nB}\,
 \frac{\beta}{\Gamma(1-\beta)}.
\]
A branch-cut integral retains the vanishing factor $\eps$ in both the
main term and its error; it yields an explicit first correction and an
arbitrary fixed-order expansion. The critical inverse has a convergent
parameter-uniform Hahn chart, whose renormalized endpoint is a
dilogarithm with its logarithmic singularity.

For the associated Poisson action process, conditioned on total weight
$n$, deleting the largest action leaves the unconditioned independent
Poisson process in total variation, for every joint limit
$n\to\infty$, $\eps\downarrow0$. When $n\eps\to\infty$, the largest
action satisfies
\[
 \frac{J-B}{\eps B}+\log\eps\ \Longrightarrow\ -Z,
 \qquad
 \E e^{itZ}=\exp\{-\tfrac\pi2|t|-it\log|t|\}.
\]
This gives a sharp Landau-distribution cutoff window and its quantiles.
When $n\eps\to\lambda<\infty$, the deficit $n-J$ instead converges in
total variation to a discrete compound-Poisson law. These limits match
as $\lambda\to\infty$. All proofs are conventional; reproducible exact
finite checks and numerical diagnostics accompany the paper.
\end{abstract}
\vfill
\noindent\textbf{Status.} This is a model-specific research manuscript with
proofs, not a claim of a general transseries breakthrough or established
publication priority. The classical analytic and condensation mechanisms
are credited. No new Lean verification is claimed.
\end{titlepage}
\tableofcontents
\clearpage

\section{The repository question and the contribution}
\label{sec:introduction}

The critical countable-action model in \repo's article
\emph{Beyond Finite-Action Folds: Critical Hahn Transseries, Stable Sector
Asymptotics, and Sharp Action Budgets} has a power-law exponent
$1<\alpha<2$. Its Question~4 asks for joint limits at both endpoints.
The editorial note at that question explicitly records that the later
endpoint papers treat $\alpha\uparrow2$, but not $\alpha\downarrow1$
\cite{repo-critical}. The present paper addresses the lower endpoint for
the exact power-law tail with an arbitrary fixed admissible finite
prefix. The repository snapshot used throughout is
\begin{center}\small\texttt{3d5973524506411392a911470b5ddc35521568ea}.\end{center}
The index and the question were read directly at that revision
\cite{repo-index,repo-critical}. This is a statement about that snapshot,
not a claim that no other unpublished work exists.

There is a genuine obstruction to substituting $\alpha=1$ into a
fixed-exponent answer. The first moment of $j^{-1-\alpha}$ diverges as
$\alpha\downarrow1$. Critical normalization therefore sends each fixed
action weight to zero. At the same time the central stable-density
constant tends to zero. A local limit theorem with an absolute error
$o(1)$ is not enough to divide by that constant, and a fixed-$\alpha$
cutoff limit need not identify the new transition window.

The principal contributions are the following linked statements.
First, a \emph{relative}, parameter-uniform coefficient theorem includes
arbitrarily small $\alpha-1$, even values much smaller than any
exponential in $n$. Second, a total-variation deletion theorem describes
the \emph{entire} remaining action process, not only the maximum. Third,
the sharp continuous cutoff window becomes a reflected, completely
asymmetric stable law of index one; a separate discrete
compound-Poisson boundary is necessary when $n(\alpha-1)$ is bounded.
A convergent critical chart, an all-order coefficient algorithm, and a
precise matching of the two probability limits complete the picture.

\subsection{Classical background and the novelty boundary}

The substitution $w=qe^U$ and Lagrange inversion are classical. The
relation of such equations to random allocations and conditioned
branching structures is treated extensively by Janson
\cite{janson}. Hankel contours and singularity transfer are established
analytic tools \cite{flajolet-odlyzko}; the exact polylogarithm expansion
used here is the standard continuation formula \cite{dlmf}.
The principle that a rare large sum is carried by one exceptional
summand, with the other summands asymptotically unaffected, is a classical
condensation mechanism; Armend\'ariz and Loulakis prove general
conditional distribution results of that kind \cite{armendariz}.
We do not claim these mechanisms as new discoveries.

What is proved here is their explicit, jointly varying realization at a
\emph{critically normalized, degenerating mean endpoint}. The proofs
retain the constants and scales that disappear in fixed-parameter
statements. In particular, no fixed-distribution big-jump theorem or
fixed-$\alpha$ local limit theorem is invoked to justify a triangular
limit outside its hypotheses. The analytical coefficient estimate and
the deletion density are established directly below.

\subsection{Scope of the answer}

The tail in this paper is exactly $j^{-2-\eps}$ after a fixed prefix.
The critical coupling is exact. We allow the prefix coefficients to have
either sign as perturbations, but the resulting action weights must be
nonnegative. We do not cover arbitrary slowly varying tails, a prefix
whose length grows with $n$, noncritical coupling windows, or unrestricted
Borel summability and resurgence. Those are separate questions. Nor is
there a nonzero finite-mean critical model obtained simply by setting
$\eps=0$ before renormalizing.

\section{Model and exact identities}
\label{sec:model}

Throughout, $\N=\{0,1,2,\ldots\}$; empty maxima and missing second maxima are zero.
Fix a real polynomial
\begin{equation}\label{eq:P}
 P(z)=\sum_{j=1}^{J_*}p_jz^j,\qquad
 p_k^{\langle\cdot\rangle}:=\sum_{j=1}^{J_*}p_jj^k,
 \qquad \kappa=P'(1)=p_1^{\langle\cdot\rangle}.
\end{equation}
The superscript in $p_k^{\langle\cdot\rangle}$ distinguishes the finite
prefix moment from the polynomial coefficient $p_k$. Assume that, for
all sufficiently small $\eps\ge0$,
\begin{equation}\label{eq:positive}
 b_{j,\eps}=j^{-2-\eps}+p_j\ge0\quad(1\le j\le J_*),
 \qquad b_{1,\eps}>0,
\end{equation}
where $p_j=0$ for $j>J_*$. Thus the tail is positive and aperiodic. Put
\begin{align}
 \alpha&=1+\eps,& \beta&=\frac1\alpha,&
 c_{\eps}&=\frac1{\zeta(1+\eps)+\kappa},\label{eq:parameters1}\\
 A_{\eps}(z)&=c_{\eps}\bigl(\Li_{2+\eps}(z)+P(z)\bigr)
             =\sum_{j\ge1}a_{j,\eps}z^j,&
 a_{j,\eps}&=c_{\eps}b_{j,\eps},\label{eq:parameters2}\\
 d_{\eps}&=A_{\eps}(1),&
 \rho_{\eps}&=e^{-d_{\eps}},&
 K_{\eps}&=c_{\eps}\Gamma(-\alpha).
 \label{eq:parameters3}
\end{align}
All these quantities are real, and $c_{\eps},K_{\eps}>0$ for small
positive $\eps$. Criticality is the exact identity
\begin{equation}\label{eq:criticality}
 A_{\eps}'(1)=\sum_{j\ge1}j a_{j,\eps}=1.
\end{equation}
The feedback solution is the unique formal series with zero constant
term satisfying
\begin{equation}\label{eq:feedback}
 U_{\eps}(q)=A_{\eps}\bigl(qe^{U_{\eps}(q)}\bigr),
 \qquad u_{n,\eps}=[q^n]U_{\eps}(q).
\end{equation}
The analytic implicit function theorem also gives convergence near zero.
The formal recursion has nonnegative coefficients, with $u_{n,\eps}>0$.
On the positive branch, $q=w e^{-A_{\eps}(w)}$ is strictly increasing
for $0<w<1$, since $wA_{\eps}'(w)<A_{\eps}'(1)=1$.
It tends to $\rho_{\eps}$ as $w\uparrow1$, which fixes the real
branch used in the local critical chart below.

\begin{proposition}[Coefficient and conditioning identities]
\label{prop:exact}
For every $n\ge1$,
\begin{equation}\label{eq:lagrange}
 u_{n,\eps}=\frac1n[z^n]e^{nA_{\eps}(z)}.
\end{equation}
Let $\mathcal N_n=(N_j)_{j\ge1}$ be independent Poisson variables of
means $\lambda_{j,n}=n a_{j,\eps}$, and set
\begin{equation}\label{eq:S}
 S_n=\sum_{j\ge1}jN_j,\qquad
 M(\mathcal N_n)=\max\{j:N_j>0\},
\end{equation}
with maximum zero for the empty process. Then
\begin{equation}\label{eq:probability}
 u_{n,\eps}=\frac{\rho_{\eps}^{-n}}n\Prob(S_n=n),
 \qquad \E S_n=n.
\end{equation}
If $u_{n,\eps}^{(m)}$ is the coefficient obtained by retaining only
actions $j\le m$, then
\begin{equation}\label{eq:budgetexact}
 \frac{u_{n,\eps}^{(m)}}{u_{n,\eps}}
 =\Prob\bigl(M(\mathcal N_n)\le m\mid S_n=n\bigr).
\end{equation}
\end{proposition}
\begin{proof}
Set $w=qe^{U_{\eps}(q)}$. Then $w=q e^{A_{\eps}(w)}$ and
$U_{\eps}=A_{\eps}(w)$. Lagrange inversion yields
\[
 [q^n]A_{\eps}(w(q))
 =\frac1n[z^{n-1}]A_{\eps}'(z)e^{nA_{\eps}(z)}
 =\frac1n[z^n]e^{nA_{\eps}(z)},
\]
where the last step differentiates the exponential. The total Poisson
intensity is $n d_{\eps}<\infty$, so the process has finitely many atoms
almost surely. Its probability generating function is
$\exp(n(A_{\eps}(z)-d_{\eps}))$, proving
\eqref{eq:probability}. The mean is \eqref{eq:criticality}.
For \eqref{eq:budgetexact}, independence of the omitted coordinates
contributes $\exp(-n\sum_{j>m}a_{j,\eps})$; multiplication by the
truncated generating function leaves exactly
$e^{-nd_{\eps}}\exp(n\sum_{j\le m}a_{j,\eps}z^j)$.
\end{proof}

An action is a part of size $j$ counted by $N_j$. We use
\emph{condensation} below in the precise sense that, under $S_n=n$, a
single exceptional action is separated from a much smaller residual
maximum. It need not carry a positive fraction of $n$ in every regime.

\section{A convergent chart through the lower endpoint}
\label{sec:chart}

Define the regular expansion coefficients
\begin{equation}\label{eq:Dk}
 D_k(\eps)=\zeta(2+\eps-k)+p_k^{\langle\cdot\rangle}
 \quad(k\ge2).
\end{equation}
They are not moments of the infinite action distribution. For example,
the actual second moment diverges, whereas $D_2(\eps)$ is finite and
may be negative. The continuation formula for the polylogarithm gives,
for $|x|<2\pi$ and $|\arg x|<\pi$,
\begin{equation}\label{eq:polylog}
 A_{\eps}(e^{-x})=d_{\eps}-x+K_{\eps}x^{\alpha}
 +c_{\eps}\sum_{k\ge2}\frac{D_k(\eps)(-x)^k}{k!}.
\end{equation}
This is an exact convergent local expansion, not merely an asymptotic
one \cite{dlmf}.

\begin{lemma}[Removable endpoint parameters]\label{lem:removable}
The functions $c_{\eps}$ and $K_{\eps}$ extend holomorphically through
$\eps=0$, with
\begin{equation}\label{eq:smallparameters}
 c_0=0,\quad c_{\eps}=\eps+O(\eps^2),\quad
 K_0=1,\quad K_{\eps}=1-(1+\kappa)\eps+O(\eps^2).
\end{equation}
On a fixed disk $|x|<r_0<2\pi$, the function
\begin{equation}\label{eq:H}
 H_{\eps}(x)=c_{\eps}\sum_{k\ge2}
                  \frac{D_k(\eps)(-1)^kx^{k-2}}{k!}
\end{equation}
is jointly holomorphic in $(\eps,x)$ near $\eps=0$, and
$H_{\eps}(x)=O(\eps)$ uniformly on smaller disks.
\end{lemma}
\begin{proof}
Use $\zeta(1+\eps)=\eps^{-1}+\gamma+O(\eps)$ and
\[
 \Gamma(-1-\eps)=\frac{\Gamma(1-\eps)}{\eps(1+\eps)},
 \qquad \Gamma(1-\eps)=1+\gamma\eps+O(\eps^2).
\]
The poles cancel in $K_{\eps}$. The only zeta pole in the full
polylogarithm expansion near this parameter is the linear term, which
has already been removed by criticality. The remaining series converges
locally uniformly; the polynomial prefix is entire after composition
with $e^{-x}$. This proves the last assertions.
\end{proof}

Write $w=e^{-x}$ and $t=\log(\rho_{\eps}/q)$, with branches continued
from positive $q<\rho_{\eps}$. The exact normal form is
\begin{equation}\label{eq:normal}
 t=K_{\eps}x^{\alpha}+x^2H_{\eps}(x),
 \qquad U_{\eps}=d_{\eps}+t-x.
\end{equation}

\begin{theorem}[Uniform convergent Hahn chart]\label{thm:chart}
There are neighborhoods of zero, independent of sufficiently small
$\eps\ge0$, and a jointly holomorphic function $V(\eps,T,R)$ such that
\begin{align}
 T&=(t/K_{\eps})^{\beta},& R&=T^{1-\eps},\label{eq:TR}\\
 x&=T V(\eps,T,T^{1-\eps}),&
 V^{\alpha}+\frac R{K_{\eps}}V^2H_{\eps}(TV)&=1.
 \label{eq:V}
\end{align}
Here $V(0,T,R)=1$, and $V-1$ is divisible by $\eps R$ in the local
holomorphic ring. In particular,
\begin{equation}\label{eq:chartfirst}
 x=T-\frac{c_{\eps}D_2(\eps)}{2\alpha K_{\eps}}
                  T^{2-\eps}
    +O\bigl(\eps |T|^{3-\eps}+\eps^2|T|^{3-2\eps}\bigr)
\end{equation}
on every sufficiently small fixed closed subsector. The full expansion
is convergent and has support contained in
\begin{equation}\label{eq:support}
 \left\{\frac{1+k+(1-\eps)m}{1+\eps}:k,m\in\N\right\}.
\end{equation}
Taylor truncations in the independent variables $(T,R)$ have a geometric
remainder on smaller polydisks, uniformly in $\eps$.
\end{theorem}
\begin{proof}
Near $V=1$ define $V^{\alpha}$ by the local logarithm. The left side of
\eqref{eq:V} minus one is jointly holomorphic and its derivative in $V$
at $(T,R,V)=(0,0,1)$ is $1+\eps$, bounded away from zero. The
parameter-dependent implicit function theorem applies on a uniform
polydisk. At $\eps=0$, $H_0=0$ and the equation has solution $V=1$.
At $R=0$ the same is true for every parameter. The Taylor expansion thus
has the factor $\eps R$. Its coefficient linear in $R$ is
$-H_{\eps}(T)/(\alpha K_{\eps})$ at $V=1$; expanding $H_{\eps}(T)$ at
zero gives \eqref{eq:chartfirst}. Substitution of $R=T^{1-\eps}$ gives
\eqref{eq:support}. Cauchy estimates on a slightly larger polydisk give
the geometric remainder. No divergent formal expansion is being summed
in this statement.
\end{proof}

The confluent coordinate can be written without hiding its large
parameter combination:
\begin{equation}\label{eq:confluent}
 T=t\exp\left(-\frac{\eps}{1+\eps}\log t
              -\frac{\log K_{\eps}}{1+\eps}\right).
\end{equation}
Expanding the exponential in powers of $\eps$ is inappropriate when
$\eps\log(1/|t|)$ is not small. Keeping it exact is the lower-endpoint
counterpart of a uniform critical chart.

\begin{proposition}[Renormalized dilogarithmic endpoint]
\label{prop:dilog}
On every compact subset of $|z|<1$,
\begin{equation}\label{eq:diloglimit}
 \frac{U_{\eps}(\rho_{\eps}z)}{\eps}
 \longrightarrow H_0^*(z):=\Li_2(z)+P(z).
\end{equation}
Its continuation near $z=e^{-t}=1$ has expansion
\begin{equation}\label{eq:diloglocal}
 H_0^*(e^{-t})=\zeta(2)+P(1)+t\log t-(1+\kappa)t
       +\sum_{k\ge2}\frac{D_k(0)(-t)^k}{k!}.
\end{equation}
Thus the nontrivial endpoint is logarithmic after division by $\eps$;
the unscaled limit is zero.
\end{proposition}
\begin{proof}
On compact subdisks, $A_{\eps}(z)/\eps\to H_0^*(z)$ locally uniformly,
while the implicit solution is $O(\eps)$. Substituting into the feedback
equation proves \eqref{eq:diloglimit}; $\rho_{\eps}\to1$ does not change
it. To obtain \eqref{eq:diloglocal}, combine the cancellation of the
linear zeta pole with the expansion of $\Gamma(-1-\eps)t^{1+\eps}$,
or differentiate \eqref{eq:normal} at $\eps=0$. The polynomial adds
$P(1)-\kappa t+\sum_{k\ge2}p_k^{\langle\cdot\rangle}(-t)^k/k!$.
\end{proof}

\section{The uniform coefficient theorem}
\label{sec:coefficient}

For $n\ge1$ define
\begin{equation}\label{eq:BJ}
 B=B_{n,\eps}=(nK_{\eps})^{1/(1+\eps)},\qquad
 h=\frac n{B^2},\qquad
 J_r(\eps)=\frac{\Gamma((r+1)/\alpha)}{\pi\alpha}
            \sin\!\left(\frac{\pi\eps(r+1)}\alpha\right).
\end{equation}
The basic constant has several equivalent forms:
\begin{equation}\label{eq:J0}
 J_0(\eps)=\frac{\beta}{\Gamma(1-\beta)}
          =-\frac1{\Gamma(-\beta)}>0,
 \qquad J_0(\eps)\sim\eps.
\end{equation}
For each fixed nonnegative integer $r$,
\begin{equation}\label{eq:Jlimit}
 J_r(\eps)\sim(r+1)!\eps,\qquad
 \frac{J_r(\eps)}{J_0(\eps)}\longrightarrow(r+1)!.
\end{equation}
Individual $J_r$ with large $r$ can have either sign away from the
endpoint; the formulas below retain their exact values.

\begin{theorem}[Relative asymptotics uniform at a vanishing amplitude]
\label{thm:uniform}
For each fixed admissible prefix $P$, there are constants
$\eps_*>0$, $C>0$, $\eta>0$, and $n_0$, with $\eps_*<1/4$, such that
for all $n\ge n_0$ and $0<\eps\le\eps_*$,
\begin{equation}\label{eq:uniform}
 u_{n,\eps}=
 \frac{\rho_{\eps}^{-n}J_0(\eps)}{nB}
 \left[1+O\left(\eps\frac n{B^2}+n^2e^{-\eta n}\right)\right].
\end{equation}
The $O$-constant is independent of $n$ and $\eps$ in this range.
Consequently the leading equivalent holds for \emph{every} joint limit
$n\to\infty$, $\eps\downarrow0$, without a lower bound on $n\eps$ or
on $\eps$ itself.
\end{theorem}

The theorem is stronger than replacing a fixed-exponent stable law by
its limit. The main probability is of order $\eps/B$, not $1/B$.
The proof in Section~\ref{sec:hankel} keeps the small sine factor in the
integral and also in its error. This is the key uniformity mechanism.

\begin{corollary}[Macroscopic crossover of sector coefficients]
\label{cor:window}
Suppose $\eps\to0$ and $\eps\log n\to\tau\in[0,\infty)$. Then
\begin{equation}\label{eq:window}
 \frac Bn\longrightarrow e^{-\tau},\qquad
 u_{n,\eps}\sim \frac{\eps e^{\tau}}{n^2}\rho_{\eps}^{-n}.
\end{equation}
When $\eps\log n\to\infty$, the general equivalent
\eqref{eq:uniform} still applies, while $B/n\to0$.
\end{corollary}
\begin{proof}
Take logarithms:
$\log(B/n)=\beta\log K_{\eps}-\eps\beta\log n$.
Use $K_{\eps}\to1$, $\beta\to1$, and $J_0/\eps\to1$.
\end{proof}

\begin{corollary}[Microscopic coefficient boundary]
\label{cor:microcoeff}
If $n\eps\to\lambda\in[0,\infty)$, then
\begin{equation}\label{eq:microcoeff}
 u_{n,\eps}\sim
 \frac{\eps}{n^2}\exp\{\lambda(\zeta(2)+P(1))\}.
\end{equation}
\end{corollary}
\begin{proof}
Now $\eps\log n\to0$, while
$nd_{\eps}=n\eps(\zeta(2)+P(1))+o(1)$.
\end{proof}

\section{Proof by a branch-cut integral}
\label{sec:hankel}

This section proves the uniform theorem without using a
parameter-dependent local limit theorem. Let
\begin{equation}\label{eq:regularR}
 R_{\eps}(r)=c_{\eps}\sum_{k\ge2}\frac{D_k(\eps)r^k}{k!}.
\end{equation}
For fixed sufficiently small $r_0>0$, uniformly in the parameter,
\begin{equation}\label{eq:Rbound}
 |R_{\eps}(r)|\le C\eps r^2\qquad(0\le r\le r_0).
\end{equation}
On the upper bank of the cut $z=e^r>1$ one has
$-\log z=-r-i0$. Equation~\eqref{eq:polylog} therefore gives
\begin{equation}\label{eq:cutexponent}
 A_{\eps}(e^r+i0)-d_{\eps}-r
 =-K_{\eps}r^{\alpha}\cos(\pi\eps)
   +iK_{\eps}r^{\alpha}\sin(\pi\eps)+R_{\eps}(r).
\end{equation}
The regular part is the same real number on both banks.

\begin{lemma}[Uniform cut representation]\label{lem:cut}
After decreasing $r_0$ and $\eps_*$ if necessary,
\begin{align}
 \Prob(S_n=n)
 &=\frac1\pi\int_0^{r_0}
 e^{-nK_{\eps}r^{\alpha}\cos(\pi\eps)+nR_{\eps}(r)}
 \sin\bigl(nK_{\eps}r^{\alpha}\sin(\pi\eps)\bigr)\dd r
 +E_{n,\eps},\label{eq:cut}\\
 |E_{n,\eps}|&\le Cn\eps e^{-\eta n}.
 \label{eq:arcerror}
\end{align}
\end{lemma}
\begin{proof}
Start with the coefficient integral for
$e^{n(A_{\eps}(z)-d_{\eps})}$. The polylogarithm continues to the plane
slit along $[1,\infty)$. Deform the circle about zero to the circle
$|z|=e^{r_0}$ together with the two banks of the cut from $1$ to
$e^{r_0}$. A small circle about $1$ contributes zero in the limit because
the integrand is bounded there. Conjugacy of the banks and
$\dd z/z=\dd r$ give \eqref{eq:cut} with the stated sign.

On the outer circle, including its two limiting endpoints on the cut,
$\Li_{2+\eps}(z)+P(z)$ is uniformly bounded for small $\eps$. Since
$c_{\eps}=O(\eps)$, decrease $\eps_*$ until the real part of
$A_{\eps}(z)-d_{\eps}-\log z$ is at most $-\eta<0$ there.
To preserve the small factor in the error, subtract the constant
$e^{-nd_{\eps}}$ inside the coefficient integrand. Its integral around
the full outer circle is zero for $n\ge1$, and it has no cut jump.
The inequality $|e^v-1|\le |v|e^{|v|}$ then bounds the remaining outer
integral by $Cn\eps e^{-\eta n}$. This subtraction is important when
$\eps$ is much smaller than $e^{-n}$.
\end{proof}

\begin{lemma}[The elementary Hankel moments]\label{lem:moments}
For $r\ge0$ an integer and small positive $\eps$,
\begin{equation}\label{eq:momentintegral}
 \frac1\pi\int_0^\infty
 s^r e^{-s^{\alpha}\cos(\pi\eps)}
 \sin\bigl(s^{\alpha}\sin(\pi\eps)\bigr)\dd s
 =J_r(\eps).
\end{equation}
Moreover, for fixed $r$, the integral of the absolute value is $O_r(\eps)$
uniformly in $\eps$.
\end{lemma}
\begin{proof}
Put $y=s^{\alpha}$ and take the imaginary part of
\[
 \frac1{\pi\alpha}\int_0^\infty
 y^{(r+1)/\alpha-1}\exp(-e^{-i\pi\eps}y)\dd y
 =\frac{\Gamma((r+1)/\alpha)}{\pi\alpha}
   e^{i\pi\eps(r+1)/\alpha}.
\]
The absolute bound follows from
$|\sin(s^{\alpha}\sin\pi\eps)|\le C\eps s^{\alpha}$ and a fixed
positive lower bound for $\cos\pi\eps$. Euler's reflection formula
then gives \eqref{eq:J0}, and Taylor expansion of the sine gives
\eqref{eq:Jlimit}.
\end{proof}

\begin{proof}[Proof of Theorem~\ref{thm:uniform}]
In \eqref{eq:cut} set $s=Br$, so $nK_{\eps}r^{\alpha}=s^{\alpha}$.
The part with $R_{\eps}=0$, extended to infinity, is $J_0(\eps)/B$.
The extension error is $O(\eps n e^{-\eta n})$, after adjusting constants,
by the last bound in Lemma~\ref{lem:moments} and
$(Br_0)^{\alpha}=nK_{\eps}r_0^{\alpha}$.

Choose $r_0$ so small that \eqref{eq:Rbound} implies
\[
 |nR_{\eps}(s/B)|\le\tfrac12\cos(\pi\eps)s^{\alpha}
 \quad(0\le s\le Br_0).
\]
This follows since $r^2/r^{\alpha}=r^{2-\alpha}$ is small uniformly on
that interval. Also
$|nR_{\eps}(s/B)|\le C\eps(n/B^2)s^2$. Hence
\begin{align*}
 &\frac1B\int_0^{Br_0}
 e^{-s^{\alpha}\cos(\pi\eps)}
 |e^{nR_{\eps}(s/B)}-1|
 \bigl|\sin(s^{\alpha}\sin\pi\eps)\bigr|\dd s\\
 &\hspace{14mm}\le
 \frac{C\eps^2}{B}\frac n{B^2}
 \int_0^\infty e^{-c s^{\alpha}}s^{\alpha+2}\dd s
 \le\frac{C\eps^2}{B}\frac n{B^2}.
\end{align*}
The integral is bounded uniformly for $1\le\alpha\le1+\eps_*$. Divide
by $J_0(\eps)/B\asymp\eps/B$. Since $B\le Cn$, the outer error divided
by the same main term is $O(n^2e^{-\eta n})$. This proves
\eqref{eq:uniform} using \eqref{eq:probability}. Finally,
$h=n/B^2\le Cn^{-(1-\eps_*)/(1+\eps_*)}\to0$ uniformly, so the stated
joint-limit conclusion follows.
\end{proof}

\begin{remark}[What the proof avoids]
The contour is a coefficient contour for the explicitly known function
$\exp(nA_{\eps})$, not a conjectured global continuation of the inverse
$U_{\eps}$. The proof therefore does not require a uniform global
inverse-function domain. It also does not divide an absolute stable
approximation by a density constant tending to zero.
\end{remark}

\section{All fixed orders and the first correction}
\label{sec:allorders}

For a finite multi-index $\nu=(\nu_2,\nu_3,\ldots)$ set
\begin{equation}\label{eq:multi}
 L(\nu)=\sum_{k\ge2}\nu_k,\quad
 r(\nu)=\sum_{k\ge2}k\nu_k,\quad
 d(\nu)=r(\nu)-L(\nu)=\sum_{k\ge2}(k-1)\nu_k.
\end{equation}
There are only finitely many multi-indices of a fixed cost $d$.

\begin{theorem}[Uniform finite-order expansion]\label{thm:allorders}
For every fixed integer $N\ge0$, uniformly in the range of
Theorem~\ref{thm:uniform},
\begin{align}
 \frac{nB\rho_{\eps}^{n}u_{n,\eps}}{J_0(\eps)}
 &=\sum_{d(\nu)\le N}
 \frac{J_{r(\nu)}(\eps)}{J_0(\eps)}
 \frac{(nc_{\eps})^{L(\nu)}}{B^{r(\nu)}}
 \prod_{k\ge2}\frac{D_k(\eps)^{\nu_k}}
 {(k!)^{\nu_k}\nu_k!}\notag\\
 &\quad+O_N\bigl(\eps h^{N+1}+n^2e^{-\eta n}\bigr).
 \label{eq:allorders}
\end{align}
The zero multi-index contributes one. The expansion is asymptotic to
arbitrary fixed order; convergence as $N\to\infty$ is not asserted.
\end{theorem}
\begin{proof}
Expand $e^{nR_{\eps}(s/B)}$ in \eqref{eq:cut}. Its monomial belonging to
$\nu$ has coefficient
\[
 \frac{(nc_{\eps})^{L(\nu)}}{B^{r(\nu)}}
 \prod_{k\ge2}\frac{D_k(\eps)^{\nu_k}}
 {(k!)^{\nu_k}\nu_k!}\ s^{r(\nu)}.
\]
Lemma~\ref{lem:moments} integrates it exactly. To justify a fixed
truncation, use the Taylor remainder for the analytic regular part on a
smaller disk and the integral remainder for the exponential. Each
nonconstant term carries at least one factor $c_{\eps}=O(\eps)$.
Since $B\le Cn$, a monomial of cost $d$ satisfies
\[
 \frac{n^{L}}{B^r}
 =h^{r-L}\left(\frac Bn\right)^{r-2L}
 \le C_N h^d\quad(d\le N+1).
\]
More explicitly, for $N\ge1$, first retain the terms $k\le N+1$
in $R_{\eps}$. Its omitted tail in the exponent has bound
$C_N\eps n B^{-N-2}s^{N+2}$, which is at most
$C_N\eps h^{N+1}s^{N+2}$. The retained exponent $Q(s)$ obeys
$|Q(s)|\le C_N\eps h s^2$ on the cut. Expanding $e^{Q(s)}$ through
power $N$ leaves an error bounded by
$C_N\eps^{N+1}h^{N+1}s^{2N+2}e^{|Q(s)|}$. Among those finitely many
powers, the monomials of cost greater than $N$ have the same required
$h^{N+1}$ factor. By making the cut radius smaller if needed, both
$e^{|Q(s)|}$ and the analogous factor for the analytic tail are absorbed
by the decaying $e^{-c s^{\alpha}}$, just as in the preceding proof.
The case $N=0$ is Theorem~\ref{thm:uniform}. The sine gives another
factor $O(\eps)$ before division by $J_0\asymp\eps$. The relative
remainder is therefore $O_N(\eps h^{N+1})$.
Polynomial factors in the exponentially small contour tails are
absorbed by decreasing $\eta$. This proves \eqref{eq:allorders}.
\end{proof}

\begin{corollary}[An explicit first correction]\label{cor:first}
Let
\begin{equation}\label{eq:delta}
 \Delta_{n,\eps}=
 \frac{nc_{\eps}D_2(\eps)}{2B^2}\,
 \frac{J_2(\eps)}{J_0(\eps)}.
\end{equation}
Then
\begin{equation}\label{eq:first}
 u_{n,\eps}=\frac{\rho_{\eps}^{-n}J_0}{nB}
 \left[1+\Delta_{n,\eps}
 +O\bigl(\eps h^2+n^2e^{-\eta n}\bigr)\right].
\end{equation}
As $\eps\downarrow0$,
\begin{equation}\label{eq:firstlimit}
 \frac{\Delta_{n,\eps}}{\eps n/B^2}
 \longrightarrow 3\left(-\frac12+p_2^{\langle\cdot\rangle}\right).
\end{equation}
In particular, the unperturbed prefix gives $-3/2$.
\end{corollary}
\begin{proof}
At cost one the only multi-index is $\nu_2=1$. Apply
$c_{\eps}/\eps\to1$, $D_2(0)=-1/2+p_2^{\langle\cdot\rangle}$, and
$J_2/J_0\to6$.
\end{proof}

For clarity, the next cost adds exactly
\begin{equation}\label{eq:secondcost}
 \frac{nc_{\eps}D_3}{6B^3}\frac{J_3}{J_0}
 +\frac{n^2c_{\eps}^2D_2^2}{8B^4}\frac{J_4}{J_0}.
\end{equation}
This is a finite, directly implementable expansion. The first prefix
moment enters the exact core $K_{\eps}$; the second enters the first
relative correction. There is no unexplained claim of prefix
independence at orders where a prefix moment is visible.

\section{Deleting the condensate: the entire cloud in total variation}
\label{sec:deletion}

Let $\mathcal N_n^*$ have the law of $\mathcal N_n$ conditioned on
$S_n=n$. Remove one occurrence of its largest action and denote the
remaining multiset by $\mathcal R_n$. The definition is unambiguous even
when the maximum is repeated: remove one atom at that size. Write $J_n$
for the removed size. All probability statements about $J_n$ henceforth
use the conditioned law.

\begin{lemma}[A deterministic first-order cloud deficit]
\label{lem:cloudshift}
For every joint limit $n\to\infty$, $\eps\downarrow0$, under the
\emph{unconditioned} law,
\begin{equation}\label{eq:cloudshift}
 \frac{S_n-n}{B}\Longrightarrow-1,
 \qquad \frac{M(\mathcal N_n)}B\longrightarrow0
 \quad\hbox{in probability}.
\end{equation}
\end{lemma}
\begin{proof}
The logarithm of the characteristic function of the centered sum is,
using \eqref{eq:polylog},
\begin{equation}\label{eq:cloudcf}
 \log\E e^{it(S_n-n)/B}
 =(-it)^{\alpha}+O_t\left(\frac{nc_{\eps}}{B^2}\right).
\end{equation}
The error tends to zero uniformly, and $(-it)^{\alpha}\to-it$, with the
principal logarithm and continuous value at $t=0$. L\'evy's continuity
theorem gives the first assertion. For each fixed $d>0$, once $dB>J_*$,
\[
 \Prob(M(\mathcal N_n)>dB)
 \le n c_{\eps}\sum_{j>dB}j^{-1-\alpha}
 \le C_d n c_{\eps}B^{-\alpha}
 =C_d\frac{c_{\eps}}{K_{\eps}}=O_d(\eps).
\]
This proves the second assertion.
\end{proof}

\begin{theorem}[Total-variation deletion theorem]\label{thm:deletion}
For every joint limit $n\to\infty$, $\eps\downarrow0$,
\begin{equation}\label{eq:TV}
 \TV\bigl(\Law(\mathcal R_n),\Law(\mathcal N_n)\bigr)\longrightarrow0.
\end{equation}
The conditioned largest action is unique with probability tending to
one, and
\begin{equation}\label{eq:Jfirst}
 \frac{J_n}{B}\longrightarrow1,
 \qquad \frac{J_n^{(2)}}B\longrightarrow0
 \quad\hbox{in probability},
\end{equation}
where $J_n^{(2)}$ is the second largest action, with multiplicities.
\end{theorem}
\begin{proof}
Let $p_n=\Prob(S_n=n)$. For an unconditioned multiset $\nu$ with total
weight $s$ and maximum $m$, define
\begin{equation}\label{eq:density}
 W_n(\nu)=\frac{\lambda_{n-s,n}}{p_n}
          \1_{\{n-s>m,\ n-s\ge1\}},
\end{equation}
where a rate with nonpositive index is zero. The Poisson deletion
identity gives, for every bounded nonnegative function $f$,
\begin{equation}\label{eq:Mecke}
 \E\bigl[f(\mathcal R_n);\ \mathcal N_n^*\text{ has a unique maximum}\bigr]
 =\E\bigl[f(\mathcal N_n)W_n(\mathcal N_n)\bigr].
\end{equation}
Indeed, adding a strictly larger atom of size $j$ to $\nu$ multiplies its
Poisson mass by $\lambda_{j,n}$; the total-weight condition forces
$j=n-s$. Alternatively, sum the usual Poisson identity
$\E[N_jf(\mathcal N_n-\delta_j)]=\lambda_{j,n}\E f(\mathcal N_n)$
with the strict-largest condition. In particular, $\E W_n\le1$.

By Lemma~\ref{lem:cloudshift}, $j=n-S_n$ satisfies $j/B\to1$ in
probability and exceeds the unconditioned maximum with probability
 tending to one. It is eventually beyond the fixed prefix, except on an
event of probability tending to zero. There,
\begin{equation}\label{eq:ratequotient}
 \frac{\lambda_{j,n}}{p_n}
 =\frac{c_{\eps}/K_{\eps}}{J_0(\eps)}
   \left(\frac Bj\right)^{1+\alpha}(1+o(1))
 \longrightarrow1,
\end{equation}
because $c_{\eps}/K_{\eps}=1/\Gamma(-1-\eps)\sim\eps$ and
$J_0\sim\eps$. The $o(1)$ in the denominator is the relative error
proved in Theorem~\ref{thm:uniform}. Hence $W_n\to1$ in probability.

This convergence automatically upgrades to $L^1$: bounded convergence
in probability gives $\E\min(W_n,1)\to1$; together with $\E W_n\le1$
it yields $\E W_n\to1$ and
\[
 \E|W_n-1|=\E W_n+1-2\E\min(W_n,1)\longrightarrow0.
\]
Equation~\eqref{eq:Mecke} now proves both that repeated maxima have
probability tending to zero and that the subprobability law on unique
maxima converges in total variation to the unconditioned law. Adding the
vanishing exceptional mass proves \eqref{eq:TV}.

Under the conditioned law, $J_n=n-S(\mathcal R_n)$ and
$J_n^{(2)}=M(\mathcal R_n)$. Total variation allows these
$n$-dependent measurable maps, so \eqref{eq:Jfirst} follows from
Lemma~\ref{lem:cloudshift}.
\end{proof}

\begin{remark}[Why the mean is not contradicted]
The unconditioned process has mean total weight exactly $n$, although
its typical total weight is $n-B+o_{\Prob}(B)$. Rare very large jumps
carry the missing mean. Total variation does not preserve expectations
of unbounded, non-uniformly-integrable observables. Thus deletion
convergence is compatible with the deterministic bound
$S(\mathcal R_n)<n$.
\end{remark}

The proof is a triangular-array version of a classical big-jump
deletion mechanism. Its short final density argument is possible only
after the vanishing central probability has been estimated
\emph{relatively}. This connects the analytic part of the paper to the
probabilistic part without an unproved uniform conditioning step.

\section{Landau fluctuations and sharp continuous budgets}
\label{sec:landau}

\begin{definition}[The endpoint stable law]\label{def:Z}
Let $Z$ be the probability law with characteristic function
\begin{equation}\label{eq:LandauCF}
 \varphi_Z(t)=\exp\{-it\log(-it)\}
 =\exp\{-\tfrac\pi2|t|-it\log|t|\},
 \quad \varphi_Z(0)=1.
\end{equation}
Let $G(z)=\Prob(-Z\le z)$ be its reflected distribution function.
\end{definition}
The logarithm in the first expression is principal. The specified
characteristic function fixes location and scale; it avoids ambiguities
among conventions for the Landau distribution. Existence also follows
from the compound-Poisson approximation in Section~\ref{sec:matching}.
Since $|\varphi_Z(t)|=e^{-\pi|t|/2}$, Fourier inversion gives a smooth
real-analytic density. It is nonnegative and not identically zero;
therefore it cannot vanish on an interval. The distribution functions of
$Z$ and $-Z$ are continuous and strictly increasing.

\begin{lemma}[The necessary microscopic scale]\label{lem:scale}
If $\eps\to0$ and $n\eps\to\infty$, then
\begin{equation}\label{eq:smallscale}
 a_n:=\eps B\to\infty,
 \qquad \frac n{\eps B^2}\to0.
\end{equation}
\end{lemma}
\begin{proof}
The bounded factors $K_{\eps}^{\beta}$ can be suppressed. If
$\eps<n^{-1/2}$, then $\eps\log n\to0$, $B\sim n$, and the two
assertions reduce to $n\eps\to\infty$. If $\eps\ge n^{-1/2}$, then
$\eps B\ge n^{1/2-o(1)}$ and
$n/(\eps B^2)\le n^{-1/2+o(1)}$. These two cases cover all indices.
\end{proof}

\begin{theorem}[Landau law for the cloud and condensate]
\label{thm:landau}
Assume $n\to\infty$, $\eps\downarrow0$, and $n\eps\to\infty$.
Under the unconditioned law,
\begin{equation}\label{eq:cloudLandau}
 Z_n:=\frac{S_n-n+B}{\eps B}-\log\eps\Longrightarrow Z.
\end{equation}
Under the conditioned law,
\begin{equation}\label{eq:condLandau}
 \frac{J_n-B}{\eps B}+\log\eps\Longrightarrow-Z.
\end{equation}
The same limiting law applies for every fixed admissible prefix, provided
the \emph{exact} prefix-dependent $K_{\eps}$ is used in $B$.
\end{theorem}
\begin{proof}
Using \eqref{eq:polylog} at frequency $t/(\eps B)$ gives
\begin{align}
 \log\E e^{itZ_n}
 &=\left(\frac{-it}{\eps}\right)^{1+\eps}
   +\frac{it}{\eps}-it\log\eps
   +O_t\left(\frac n{\eps B^2}\right).
 \label{eq:LandauExpand}
\end{align}
The remainder tends to zero by Lemma~\ref{lem:scale}. For each fixed
$t\ne0$, write $L=\log(-it)-\log\eps$. Then
\[
 \left(\frac{-it}{\eps}\right)^{1+\eps}
 =\frac{-it}{\eps}e^{\eps L}
 =\frac{-it}{\eps}-itL
   +O_t\bigl(\eps(1+|\log\eps|)^2\bigr).
\]
Substitution into \eqref{eq:LandauExpand} leaves
$-it\log(-it)+o(1)$. Continuity at zero and L\'evy's theorem prove
\eqref{eq:cloudLandau}. For the conditional statement use
$J_n=n-S(\mathcal R_n)$ and Theorem~\ref{thm:deletion}. The identity of
the normalized variables has exactly the minus sign displayed in
\eqref{eq:condLandau}.
\end{proof}

The logarithmic displacement is essential. The leading location is
$B-\eps B\log\eps$, not simply $B$, on the fluctuation scale $\eps B$.
Because $-\log\eps\to\infty$, dropping it changes every fixed cutoff
quantile.

\begin{theorem}[Sharp cutoff profile and quantiles]\label{thm:budget}
Under the hypotheses of Theorem~\ref{thm:landau}, put
\begin{equation}\label{eq:budgetwindow}
 m_n(z)=\left\lfloor B+\eps B(z-\log\eps)\right\rfloor.
\end{equation}
For every real $z$,
\begin{equation}\label{eq:budgetG}
 \frac{u_{n,\eps}^{(m_n(z))}}{u_{n,\eps}}\longrightarrow G(z).
\end{equation}
For $\theta\in(0,1)$ let $q_\theta=G^{-1}(\theta)$ and define the
minimal action budget
\begin{equation}\label{eq:minbudget}
 m^*_{n,\eps}(\theta)=\min\{m\in\N:
              u_{n,\eps}^{(m)}/u_{n,\eps}\ge\theta\}.
\end{equation}
Then
\begin{equation}\label{eq:quantile}
 m^*_{n,\eps}(\theta)
 =B+\eps B(q_\theta-\log\eps)+o(\eps B).
\end{equation}
\end{theorem}
\begin{proof}
Equation~\eqref{eq:budgetexact}, Theorem~\ref{thm:landau}, and
$\eps B\to\infty$ give \eqref{eq:budgetG}; the floor is negligible.
Continuity and strict increase of $G$ imply quantile convergence: for
any fixed $d>0$, the retained fraction at $m_n(q_\theta-d)$ is eventually
less than $\theta$, and at $m_n(q_\theta+d)$ it is eventually greater.
Monotonicity sandwiches the minimal budget, and then $d\downarrow0$.
\end{proof}

\begin{corollary}[Coarse cutoff step]\label{cor:step}
For fixed $0<s<1$ the retained fraction at $m=\lfloor sB\rfloor$ tends to
zero, and for fixed $s>1$ it tends to one. At $m=\lfloor B\rfloor$ it
also tends to zero under $n\eps\to\infty$.
\end{corollary}
\begin{proof}
The strict inequalities follow from $J_n/B\to1$. At $m=B$, the
normalized threshold in \eqref{eq:condLandau} is $\log\eps+o(1)\to-\infty$;
tightness proves the final assertion.
\end{proof}

\subsection{The residual extremes}

\begin{theorem}[Poisson residual extremes]\label{thm:extremes}
Under $n\eps\to\infty$, the point process of residual action sizes,
rescaled by $a_n=\eps B$, converges on $(0,\infty)$, vaguely on sets
bounded away from zero, to a Poisson point process of intensity
$x^{-2}\dd x$. In particular,
\begin{equation}\label{eq:secondlargest}
 \Prob\left(\frac{J_n^{(2)}}{\eps B}\le x\right)
 \longrightarrow e^{-1/x}\qquad(x>0).
\end{equation}
No independence between this extreme process and the Landau cloud sum
is claimed.
\end{theorem}
\begin{proof}
For the unconditioned process the rescaled intensity on a compact
interval in $(0,\infty)$ is a Riemann sum with factor
\[
 n c_{\eps}a_n^{-\alpha}
 =\frac{c_{\eps}}{K_{\eps}}\eps^{-\alpha}
 \longrightarrow1,
\]
since $c_{\eps}/K_{\eps}\sim\eps$ and
$\eps^{-\eps}\to1$. The lattice spacing $1/a_n$ tends to zero.
Independence of disjoint Poisson counts gives convergence of Laplace
functionals. The finite prefix disappears from every such interval.
For the maximum, the exact omitted intensity satisfies
$n\sum_{j>a_nx}a_{j,\eps}\to1/x$, so its zero-count probability tends
to $e^{-1/x}$. Transfer both conclusions to the residual process using
Theorem~\ref{thm:deletion}.
\end{proof}

\section{The discrete compound-Poisson boundary}
\label{sec:poisson}

Let $\mathcal C_\lambda$ be the independent Poisson action process with
rates $\lambda b_{j,0}=\lambda(j^{-2}+p_j)$, and let $C_\lambda$ be its
total weight. Its total intensity is finite even though its mean is
infinite when $\lambda>0$. Its generating function is
\begin{equation}\label{eq:CPpgf}
 \E z^{C_\lambda}
 =\exp\left\{\lambda\bigl(\Li_2(z)+P(z)-\zeta(2)-P(1)\bigr)\right\}.
\end{equation}
For $\lambda=0$ this is the point mass at zero.

\begin{theorem}[Microscopic cloud and deficit in total variation]
\label{thm:CP}
If $n\eps\to\lambda\in[0,\infty)$, then
\begin{equation}\label{eq:CPtv}
 \TV\bigl(\Law(\mathcal R_n),\Law(\mathcal C_\lambda)\bigr)\to0,
 \qquad
 \TV\bigl(\Law(n-J_n),\Law(C_\lambda)\bigr)\to0.
\end{equation}
Consequently, for every fixed integer $k\ge0$,
\begin{equation}\label{eq:CPbudget}
 \frac{u_{n,\eps}^{(n-k)}}{u_{n,\eps}}
 \longrightarrow\Prob(C_\lambda\ge k).
\end{equation}
\end{theorem}
\begin{proof}
For each $j$, $n a_{j,\eps}\to\lambda b_{j,0}$. Since $n c_{\eps}$
is bounded, dominated convergence with the summable tail $j^{-2}$ gives
\[
 \sum_{j\ge1}|n a_{j,\eps}-\lambda b_{j,0}|\longrightarrow0.
\]
Couple independent Poisson counts by their common intensities and
independent excess counts. The probability of any disagreement is at
most the sum above, so the unconditioned processes converge in total
variation. Combine this with Theorem~\ref{thm:deletion}. The total-weight
map preserves the total-variation bound, proving the deficit claim.
Finally $J_n\le n-k$ is exactly $n-J_n\ge k$.
\end{proof}

For example, the limiting retained fraction when the single largest
possible action is omitted is
\begin{equation}\label{eq:lastaction}
 \frac{u_{n,\eps}^{(n-1)}}{u_{n,\eps}}
 \longrightarrow1-e^{-\lambda(\zeta(2)+P(1))}.
\end{equation}
When $n\eps\to0$, essentially all the coefficient comes from
configurations whose largest action is exactly $n$; every cutoff below
$n$ loses asymptotically all of it. This is not a continuous cutoff
profile. For a prescribed retained fraction at a lattice boundary,
quantile statements must respect atoms: one may compare the discrete
tails in \eqref{eq:CPbudget}, but cannot assert convergence of the
minimal integer budget at a threshold equal to a limiting tail
probability without additional information.

\section{Matching the two endpoint descriptions}
\label{sec:matching}

\begin{proposition}[Compound-Poisson to Landau matching]
\label{prop:matching}
As $\lambda\to\infty$,
\begin{equation}\label{eq:matching}
 \frac{C_\lambda}{\lambda}-\log\lambda-(1+\kappa)
 \Longrightarrow Z.
\end{equation}
\end{proposition}
\begin{proof}
Put $x=-it/\lambda$ in \eqref{eq:diloglocal}. The logarithm of the
characteristic function of $C_\lambda/\lambda$ is
\[
 \lambda\bigl(H_0^*(e^{it/\lambda})-H_0^*(1)\bigr)
 =-it\log(-it)+it\log\lambda+it(1+\kappa)+O_t(\lambda^{-1}).
\]
Subtracting the two centering terms gives \eqref{eq:LandauCF} in the
limit. L\'evy's theorem proves the claim and, independently, establishes
existence of the probability law in Definition~\ref{def:Z}.
\end{proof}

This calculation checks the otherwise delicate constant in the
logarithmic displacement. For example, in the matching overlap
\begin{equation}\label{eq:matchingoverlap}
 \eps\log^2 n\longrightarrow0,\qquad n\eps\longrightarrow\infty,
\end{equation}
write $\lambda_n=n\eps$. Expansion of the exact core gives
\begin{equation}\label{eq:Bexpand}
 B=n-n\eps(\log n+1+\kappa)
       +O\bigl(n\eps^2(1+\log n)^2\bigr).
\end{equation}
It follows that the conditional deficit, on the scale $\lambda_n$,
is centered at
\begin{equation}\label{eq:centerMatch}
 n-B+\eps B\log\eps
 =\lambda_n\bigl(\log\lambda_n+1+\kappa\bigr)+o(\lambda_n).
\end{equation}
Together with Theorem~\ref{thm:landau} this gives the same centered
Landau law as \eqref{eq:matching}. Condition
\eqref{eq:matchingoverlap} is stated explicitly: without it, replacing
the exact $B$ by its first-order expansion need not be accurate on the
fluctuation scale. The exact-core formulation requires no such extra
condition.

\subsection{A regime map}

The two combinations $\eps\log n$ and $n\eps$ play different roles.
The first controls the fraction of total mass in the exceptional action;
the second determines whether the residual cloud is continuous after
rescaling or remains discrete. Table~\ref{tab:regimes} summarizes this.
All entries assume $\eps\downarrow0$ and $n\to\infty$.

\begin{table}[htbp]\centering\small
\begin{tabular}{@{}p{.25\textwidth}p{.28\textwidth}p{.38\textwidth}@{}}
\toprule
Parameter regime & Largest action & Residual description and budget \\
\midrule
$\eps\log n\to\tau\in(0,\infty)$
& $J_n/n\to e^{-\tau}$
& Landau window of width $\eps B$, centered at $B-\eps B\log\eps$.\\[5pt]
$\eps\log n\to\infty$
& $J_n\sim B=o(n)$
& One action is still separated from the residual maximum by a factor
of order $1/\eps$; the same Landau window applies.\\[5pt]
$\eps\log n\to0$, $n\eps\to\infty$
& $J_n/n\to1$
& A diverging residual cloud; retain the exact $B$ for fluctuation-scale
accuracy.\\[5pt]
$n\eps\to\lambda\in(0,\infty)$
& $J_n=n-O_{\Prob}(1)$
& Deficit tends to $C_\lambda$; cutoff transitions are discrete.\\[5pt]
$n\eps\to0$
& $\Prob(J_n=n)\to1$
& Omitting action $n$ loses asymptotically the whole coefficient.\\
\bottomrule
\end{tabular}
\caption{The lower-endpoint regimes. ``$J_n\sim B$'' here means
$J_n/B\to1$ in probability, not convergence of expectations.}
\label{tab:regimes}
\end{table}

The leading coefficient remains the single expression
$\rho_{\eps}^{-n}J_0/(nB)$ throughout this table. What changes is the
structure of conditioning and the resolution needed for the action
budget. This is why a coefficient equivalent alone would not answer the
repository's cutoff question.

\section{Reproducible checks and numerical evidence}
\label{sec:computation}

The package contains \texttt{verify.py} and its recorded output
\texttt{verification\_results.json}. The program performs three distinct
kinds of checks. First, 385 exact rational assertions compare a direct
feedback recursion with Lagrange coefficients, enumerate finite action
multisets, and check the Poisson deletion and cutoff identities.
Second, an independent 80-digit recurrence at $n=100$, $\eps=0.05$
agrees with the log-space numerical implementation to a relative
difference of approximately $2.19\times10^{-15}$. Third, it evaluates
the asymptotic and probability formulas at finite parameters. These last
calculations are floating-point diagnostics, not proofs and not interval
certificates.

\subsection{Coefficient recurrence and normalization}

For a fixed target index $n$, let
$p_k=\Prob(S_n=k)$. Then
\begin{equation}\label{eq:recurrence}
 p_0=e^{-nd_{\eps}},\qquad
 p_k=\frac nk\sum_{j=1}^k j a_{j,\eps}p_{k-j}\quad(1\le k\le n).
\end{equation}
The implementation uses log-sum-exp, so it does not require an extended
floating-point exponent range. For a cutoff $m$, retain $j\le m$ in the
sum but \emph{keep the same full} $p_0=e^{-nd_{\eps}}$. The resulting
coefficient is $\Prob(S_n=k, N_j=0\text{ for }j>m)$, not the normalized
probability under a different truncated process. This distinction is
necessary for \eqref{eq:budgetexact}.

Table~\ref{tab:coefficients} compares the leading and corrected formulas
for the unperturbed prefix. All values in the tables below are generated
from the recorded JSON output; no qualitative convergence assertion is
inferred from a few finite ratios.

\begin{table}[htbp]\centering\small
\begin{tabular}{@{}rrrrr@{}}
\toprule
$\eps$ & $n$ & \shortstack{Recurrence /\\leading} & \shortstack{Recurrence /\\corrected}
& \shortstack{Observed correction /\\predicted} \\
\midrule
\inputtable{coefficient_table.tex}
\bottomrule
\end{tabular}
\caption{The first correction materially improves the coefficient
approximation. ``Exact recurrence'' means the exact recurrence evaluated
in floating point, not exact real arithmetic.}
\label{tab:coefficients}
\end{table}

The program also checks the prefix
$P(z)=z/2-z^2/16$. Here $\kappa=3/8$,
$p_2^{\langle\cdot\rangle}=1/4$, and the limiting first-correction
constant is $-3/4$, rather than $-3/2$. This checks the explicit prefix
sensitivity rather than suppressing it.

\subsection{Landau profile and finite-parameter limitations}

Fourier inversion gives a convenient formula for the distribution of
$Z$:
\begin{equation}\label{eq:CDF}
 F_Z(x)=\frac12+\frac1\pi\int_0^\infty
 e^{-\pi t/2}\frac{\sin(t(x+\log t))}{t}\dd t,
 \qquad G(z)=1-F_Z(-z).
\end{equation}
The integrable logarithmic singularity at zero is handled by splitting
the integral. The exponential tail permits a controlled tail estimate,
although the supplied quadrature itself does not use directed rounding.
The numerical quantiles are
\begin{equation}\label{eq:quantilenumbers}
 q_{0.1}\approx-11.6492846845,\qquad
 q_{0.5}\approx-1.3557804210,\qquad
 q_{0.9}\approx1.0922545281.
\end{equation}
Table~\ref{tab:landau} deliberately shows that a sharp limit theorem need
not give a high-accuracy finite-$\eps$ profile. The errors are still
visible at the displayed parameters. No rate for
Theorem~\ref{thm:budget} is claimed.

\begin{table}[htbp]\centering\small
\begin{tabular}{@{}rrrrrr@{}}
\toprule
$n$ & $\eps$ & $z$ & cutoff $m_n(z)$ & retained fraction & $G(z)$\\
\midrule
\inputtable{landau_table.tex}
\bottomrule
\end{tabular}
\caption{Finite-parameter cutoff diagnostics. The discrepancies, rather
than only favorable ratios, are retained in the report.}
\label{tab:landau}
\end{table}

\subsection{The discrete boundary}

For $k<n/2$, the event $J_n=n-k$ has a unique largest atom, so its
probability is exactly
\begin{equation}\label{eq:deficitExact}
 \Prob(n-J_n=k)
 =\frac{n a_{n-k,\eps}\Prob(S_n=k)}{\Prob(S_n=n)}.
\end{equation}
This allows a check of the finite-$\lambda$ cloud without Monte Carlo
conditioning on a rare event. For $P=0$, $\lambda=1/2$, the mass at zero
should tend to $e^{-\pi^2/12}\approx0.4393464341$. The computed values
are $0.4343650385$ at $n=512$ and $0.4377984669$ at $n=2048$.
The recorded output includes both prefixes, $\lambda=1/2,2$, and
errors for deficits $0$ through $15$.

\subsection{What has and has not been verified}

The 385 rational assertions verify finite algebraic identities only.
The numerical recurrences check the stated formulas and normalizations
at their input parameters. The asymptotic theorems, total-variation
convergence, and parameter-uniform errors are supported by the proofs in
this paper, not by those tests. No interval-certified action budget has
been produced, and no Lean file is included. The supplied sources are
sufficient to rebuild the article and rerun every reported computation.

\section{A formalization and implementation route}
\label{sec:formalization}

The repository separates the transseries foundations from its research
reports \cite{repo-index}. Nothing in the present package should be
marked as a Lean theorem merely because related foundational files
exist. A practical formalization can nevertheless be decomposed into
small interfaces.

The finite algebra comes first: formal coefficient extraction,
Lagrange inversion in the form \eqref{eq:lagrange}, the exponential
coefficient recurrence, and the finite multiset cutoff identity. The
next layer can treat a countable summable sequence of Poisson rates and
prove the exact deletion density \eqref{eq:Mecke}. The upgrade from
$W_n\to1$ in probability and $\E W_n\le1$ to $L^1$ convergence is a
short, reusable lemma independent of transseries.

The analytic layer needs the polylogarithm continuation formula with
uniform parameter estimates on a compact slit contour, followed by the
cut deformation, the elementary gamma integral, and the uniform
exponential remainder bound. It is useful to state these as an abstract
``regular part plus one branch term'' coefficient theorem before
specializing to the zeta coefficients. Characteristic-function
convergence and the exact-core expansion then provide the probability
limits. Finally, quantile convergence follows from a general monotone
CDF lemma.

The code accompanying this paper is an implementation of finite
recurrences and diagnostics. It is not a placeholder formal proof, and
it does not replace any of those analytical obligations.

\section{Further research questions}
\label{sec:future}

\paragraph{1. Quantitative deletion and Landau approximation.}
The total-variation theorem here is qualitative. Derive an explicit
bound for the deletion distance and for the cutoff CDF error, uniformly
in both $n\eps$ and $\eps\log n$. The characteristic-function proof
suggests contributions of order
$\eps\log^2(1/\eps)$ and $n/(\eps B^2)$, but smoothing and the deletion
density introduce additional terms. These suggestions are not error
bounds proved in this paper. Such a result would explain and improve
the visible finite-parameter discrepancies in Table~\ref{tab:landau}.

\paragraph{2. An interval-certified minimal action budget.}
Combine the exact positive recurrence with rigorous logarithmic
arithmetic, certified zeta and gamma values, and a proved finite-$n$
CDF bound. The goal is a program that returns the smallest cutoff with
a guaranteed retained fraction, or a narrow certified bracket. The
asymptotic quantile formula alone is not such a program.

\paragraph{3. Slowly varying and nonexact tails.}
Replace the exact tail by
$j^{-2-\eps}L_{\eps}(j)$. Determine verifiable uniform regularity
conditions under which the deletion theorem and Landau window persist,
and identify the new centering constants. Without an analytic
continuation of the generating function, the current cut proof is not
available; a genuinely uniform local big-jump estimate would be needed.

\paragraph{4. Growing or parameter-singular finite cores.}
The proof allows any fixed admissible prefix, but not a prefix whose
length tends to infinity or whose moments diverge with $\eps$.
Find sharp criteria separating harmless prefix changes from a change
in the limiting cloud. In particular, decide when the analytic regular
part remains negligible on the scale $\eps B$.

\paragraph{5. Expanding noncritical coupling windows.}
For $U=\theta A_{\eps}(qe^U)$, criticality is $\theta=1$.
Study the joint variables $1-\theta$, $\eps$, and $n$.
A deterministic shift of the cloud competes with the exceptional action
and can move the cutoff center by its entire width. The present exact
critical normalization does not settle that crossover.

\paragraph{6. Uniform matching beyond the first Landau order.}
Construct a single approximation that retains the discrete
compound-Poisson law for bounded $n\eps$, has a computable first
Landau correction for diverging $n\eps$, and uses the exact core for
arbitrary $\eps\log n$. This would refine both boundary descriptions
without imposing \eqref{eq:matchingoverlap}.

\paragraph{7. Large-order behavior of the endpoint expansion.}
The convergent Hahn chart and the fixed-order sector expansion are
not the same kind of series. Determine the growth of the coefficients
in \eqref{eq:allorders}, an optimal truncation rule, and any additional
exponentially small sectors. No Borel or resurgence conclusion follows
from fixed-order uniformity alone.

\paragraph{8. The unnormalized exponent-one model.}
At exactly $\alpha=1$, the unnormalized kernel
$a\Li_2(z)+P(z)$ has infinite first moment, and the normalization used
here collapses to zero. Analyze its interior critical point as the
coupling tends to zero, and match that problem to the renormalized
endpoint studied here. It is a different order of limits, not a missing
substitution in \eqref{eq:uniform}.

\paragraph{9. Other analytic branch kernels and multivariate actions.}
Abstract the sine-factor argument to kernels with a branch exponent
approaching an integer and a simultaneously vanishing normalization.
Determine which multivariate constraints admit an analogue of the
one-atom deletion density, and when several exceptional actions rather
than one are forced by geometry.

\section{Conclusion}

The lower critical endpoint can be analyzed without a formal
substitution into a fixed-exponent theorem. The relevant coefficient
probability vanishes like $\eps/B$, and a branch-cut proof preserves
that small factor uniformly. Once the relative estimate is available,
a direct deletion density proves convergence of the entire conditioned
remainder to the unconditioned Poisson cloud. Its fluctuations produce
a Landau cutoff window when $n\eps$ diverges and a discrete
compound-Poisson deficit when it does not.

These results answer the lower-endpoint chart, coefficient, and cutoff
questions for the exact power-law model with a fixed admissible prefix.
They also identify precisely what remains: quantitative conditional
errors, certified budgets, nonexact tails, and noncritical or growing-core
extensions. The distinction between a proved model-specific continuation
and a general transseries theorem is maintained throughout.

\appendix
\section{Endpoint constants and sign checks}
\label{app:constants}

The expansions required to audit the centering are
\begin{align}
 c_{\eps}&=\eps-(\gamma+\kappa)\eps^2+O(\eps^3),\\
 K_{\eps}&=1-(1+\kappa)\eps+O(\eps^2),\\
 J_0(\eps)&=\eps+(\gamma-2)\eps^2+O(\eps^3),\\
 K_{\eps}^{-\beta}J_0(\eps)
 &=\eps\bigl[1+(\gamma-1+\kappa)\eps+O(\eps^2)\bigr].
\end{align}
The sign of the stable variable is checked directly from
\[
 -it\log(-it)=-\tfrac\pi2|t|-it\log|t|.
\]
For the largest action the variable is reflected because the exact
identity is $J_n=n-S(\mathcal R_n)$. Its quantile center therefore is
$B+\eps B(q_\theta-\log\eps)$, with a plus sign before $q_\theta$.
The residual total weight uses the opposite sign.

There is also a finite-$n$ perturbative check for $P=0$. At first order
in $\eps$, $u_{n,\eps}=\eps/n^2+O_n(\eps^2)$. At second order,
\begin{equation}\label{eq:perturbcheck}
 \frac{n^2u_{n,\eps}}{\eps}
 =1+\eps\left(nH_{n-1}^{(2)}+2H_{n-1}-\gamma-\log n\right)
 +O_n(\eps^2),
\end{equation}
where $H_m=\sum_{j\le m}1/j$ and
$H_m^{(2)}=\sum_{j\le m}1/j^2$. Indeed,
\[
 [z^n]\Li_2(z)^2
 =\frac{2H_{n-1}^{(2)}}{n^2}+\frac{4H_{n-1}}{n^3},
\]
which follows by partial fractions in $j$ and $n-j$.
Expanding the harmonic sums gives the relative correction
$-3\eps/(2n)$ at the endpoint, agreeing with
\eqref{eq:firstlimit}. The $O_n$ notation here is explicitly for fixed
$n$; this perturbative check is not used as a uniform proof.

\section{Rebuilding and source provenance}
\label{app:reproduce}

The archive contains the article source and compiled PDF, generated
LaTeX tables, the verification program, its recorded JSON results,
\texttt{README.md}, \texttt{requirements.txt}, and a build script.
The computational commands are
\begin{lstlisting}
python verify.py --output verification_results.json
python make_tables.py
latexmk -pdf -interaction=nonstopmode -halt-on-error article.tex
\end{lstlisting}
The \texttt{--quick} verification option omits the larger cutoff
calculations, and should write to a different output filename when the
recorded full output is to be preserved. No network access is required
by the verification program. The numerical dependencies are NumPy,
SciPy, and mpmath; exact assertions use Python's standard-library
\texttt{fractions} module.

The repository question and model were read from the pinned
\path{Critical_Hahn_Transseries_Beyond_Finite_Action_Folds/article.tex}
under \path{Analysis/Transseries/docs/series-and-transseries/}.
The motivating question is Section~``Further research questions'',
Question~4; the exact coefficient and Poisson identities appear in its
model section. The fixed-exponent chart motivates
Section~\ref{sec:chart}, but all uniform endpoint, deletion, and cutoff
arguments used here are proved in this manuscript. No assumption that
the unreviewed predecessor is correct is needed to validate these proofs.

\begin{thebibliography}{9}
\raggedright
\bibitem{repo-critical}
\repo\ repository, maintained by V.~Reshetnikov,
\emph{Beyond Finite-Action Folds: Critical Hahn Transseries, Stable Sector
Asymptotics, and Sharp Action Budgets}, editorial revision of
29 September 2026. Repository revision
\texttt{3d5973524506411392a911470b5ddc35521568ea},
\path{Analysis/Transseries/docs/series-and-transseries/Critical_Hahn_Transseries_Beyond_Finite_Action_Folds/article.tex}.
\url{https://github.com/VladimirReshetnikov/ProveIt}.

\bibitem{repo-index}
\repo\ transseries research index, maintained by V.~Reshetnikov, same revision,
\path{Analysis/Transseries/docs/series-and-transseries/README.md}.
\url{https://github.com/VladimirReshetnikov/ProveIt}.

\bibitem{dlmf}
NIST Digital Library of Mathematical Functions,
Section~25.12, especially Eq.~25.12.12 (polylogarithm continuation).
\url{https://dlmf.nist.gov/25.12}.
Accessed 29 September 2026.

\bibitem{flajolet-odlyzko}
P.~Flajolet and A.~M.~Odlyzko,
Singularity analysis of generating functions,
\emph{SIAM Journal on Discrete Mathematics} \textbf{3} (1990), 216--240.
\url{https://doi.org/10.1137/0403019}.

\bibitem{janson}
S.~Janson,
Simply generated trees, conditioned Galton--Watson trees, random
allocations and condensation,
\emph{Probability Surveys} \textbf{9} (2012), 103--252.
\url{https://arxiv.org/abs/1112.0510}.

\bibitem{armendariz}
I.~Armend\'ariz and M.~Loulakis,
Conditional distribution of heavy tailed random variables on large
deviations of their sum,
\emph{Stochastic Processes and their Applications} \textbf{121} (2011),
1138--1147.
\url{https://arxiv.org/abs/0912.1516}.
\end{thebibliography}
\end{document}
